% ─────────────────────────────────────────────────────────────
%  PICK & SERVE  |  White Paper – Entrega 2
%  Ingeniería de Software II  |  ITBA  |  Junio 2026
% ─────────────────────────────────────────────────────────────
\documentclass[12pt, a4paper]{article}

% ── Paquetes ──────────────────────────────────────────────────
\usepackage[utf8]{inputenc}
\usepackage[T1]{fontenc}
\usepackage[spanish]{babel}
\usepackage[left=2.5cm, right=2.5cm, top=3cm, bottom=3cm]{geometry}
\usepackage{graphicx}
\usepackage{float}
\usepackage{hyperref}
\usepackage{xcolor}
\usepackage{listings}
\usepackage{enumitem}
\usepackage{booktabs}
\usepackage{array}
\usepackage{caption}
\usepackage{subcaption}
\usepackage{titlesec}
\usepackage{fancyhdr}
\usepackage{parskip}
\usepackage{microtype}
\usepackage{lmodern}
\usepackage{tikz}
\usetikzlibrary{shapes.geometric, arrows.meta, positioning, fit, backgrounds}

% ── Colores de marca ──────────────────────────────────────────
\definecolor{brand}{HTML}{1A5276}
\definecolor{accent}{HTML}{2E86C1}
\definecolor{light}{HTML}{EBF5FB}
\definecolor{codebg}{HTML}{F4F6F7}
\definecolor{gray60}{HTML}{666666}

% ── Hyperref ──────────────────────────────────────────────────
\hypersetup{
  colorlinks=true,
  linkcolor=brand,
  urlcolor=accent,
  citecolor=brand
}

% ── Listings (código) ─────────────────────────────────────────
\lstset{
  backgroundcolor=\color{codebg},
  basicstyle=\ttfamily\footnotesize,
  breaklines=true,
  frame=single,
  framerule=0.4pt,
  rulecolor=\color{gray60},
  numbers=none,
  captionpos=b,
  literate={á}{{\'a}}1 {é}{{\'e}}1 {í}{{\'i}}1 {ó}{{\'o}}1 {ú}{{\'u}}1
           {Á}{{\'A}}1 {É}{{\'E}}1 {Í}{{\'I}}1 {Ó}{{\'O}}1 {Ú}{{\'U}}1
           {ñ}{{\~n}}1 {Ñ}{{\~N}}1
}

% ── Encabezado / pie ──────────────────────────────────────────
\pagestyle{fancy}
\fancyhf{}
\renewcommand{\headrulewidth}{0.4pt}
\fancyhead[L]{\small\textcolor{gray60}{Pick \& Serve — White Paper}}
\fancyhead[R]{\small\textcolor{gray60}{Ingeniería de Software II · ITBA}}
\fancyfoot[C]{\small\thepage}

% ── Secciones ─────────────────────────────────────────────────
\titleformat{\section}{\large\bfseries\color{brand}}{}{0em}{}[\titlerule]
\titleformat{\subsection}{\normalsize\bfseries\color{accent}}{}{0em}{}

% ─────────────────────────────────────────────────────────────
\begin{document}

% ── Portada ───────────────────────────────────────────────────
\thispagestyle{empty}
\begin{center}
  \vspace*{1.5cm}
  {\fontsize{28}{34}\selectfont\bfseries\color{brand} Pick \& Serve}\\[0.4cm]
  {\large\color{gray60} Liga de Predicciones de Tenis ATP}\\[2cm]

  \begin{tikzpicture}
    \fill[brand!10, rounded corners=8pt] (0,0) rectangle (12,0.08);
  \end{tikzpicture}

  \vspace{1.2cm}
  {\LARGE\bfseries White Paper}\\[0.3cm]
  {\large Arquitectura Orientada a Eventos — Entrega Final}\\[2.5cm]

  \begin{tabular}{rl}
    \textbf{Integrantes:} & Matías Mutz\\[0.2cm]
    \textbf{Materia:}     & Ingeniería de Software II\\
    \textbf{Institución:} & Instituto Tecnológico de Buenos Aires (ITBA)\\
    \textbf{Fecha:}       & Junio 2026\\
  \end{tabular}
\end{center}

\newpage
\thispagestyle{empty}
\tableofcontents
\newpage

% ─────────────────────────────────────────────────────────────
\section{Introducción y Contexto del Problema}

Organizar una liga de predicciones deportivas parece, a primera vista, un problema simple: cada usuario elige un ganador antes de cada partido y, cuando el resultado está disponible, el sistema le suma puntos. La realidad operativa, sin embargo, dista bastante de eso.

Cuando un administrador carga el resultado de un partido, se desencadena una secuencia de operaciones que están fuertemente relacionadas entre sí pero que deberían poder procesarse de forma independiente: calcular los puntos de cada predicción, actualizar el ranking global, notificar a cada usuario del resultado de su apuesta y, de ser el caso, avanzar automáticamente el estado de la ronda. Si estas operaciones se implementan de forma síncrona dentro de la misma solicitud HTTP, el sistema se vuelve frágil, lento y difícil de escalar. Un error en el envío de notificaciones, por ejemplo, no debería impedir que el ranking se actualice.

El problema central que aborda Pick \& Serve no es la liga en sí misma, sino el modelo de procesamiento subyacente: ¿cómo desacoplar operaciones relacionadas causalmente sin sacrificar consistencia, y cómo hacerlo de una manera que el sistema pueda crecer en complejidad sin volverse inmanejable?

% ─────────────────────────────────────────────────────────────
\section{Solución Propuesta: Arquitectura Orientada a Eventos}

La respuesta que propone Pick \& Serve es reemplazar el procesamiento síncrono encadenado por un modelo de publicación/suscripción (\textit{pub/sub}) mediado por un broker de mensajes. Cada acción significativa en el sistema produce un evento; los distintos consumidores reaccionan a ese evento de forma independiente, sin conocer ni depender unos de otros.

\begin{figure}[H]
\centering
\begin{tikzpicture}[
  node distance=0.9cm and 1.5cm,
  box/.style={rectangle, rounded corners=5pt, draw=brand, fill=brand!8,
              text centered, minimum height=1cm, minimum width=2.8cm,
              font=\small},
  worker/.style={rectangle, rounded corners=5pt, draw=accent, fill=accent!10,
                 text centered, minimum height=1cm, minimum width=2.8cm,
                 font=\small},
  queue/.style={cylinder, shape border rotate=90, draw=gray60!70, fill=codebg,
                minimum height=0.8cm, minimum width=1.4cm, font=\tiny,
                aspect=0.4},
  arr/.style={-Stealth, thick, color=brand},
  arrg/.style={-Stealth, thick, color=gray60}
]

%% Columna izquierda: API
\node[box] (api) {Backend API\\{\tiny FastAPI}};

%% Exchange
\node[box, right=2.2cm of api, fill=accent!15, draw=accent] (exch)
      {Exchange\\{\tiny pickandserve.events}};

%% Queues
\node[queue, right=1.8cm of exch, yshift=1.8cm]  (q1) {scoring.queue};
\node[queue, right=1.8cm of exch, yshift=0cm]     (q2) {notifications.queue};
\node[queue, right=1.8cm of exch, yshift=-1.8cm]  (q3) {notifications\\round.queue};

%% Workers
\node[worker, right=1.6cm of q1] (w1) {Scoring\\Worker};
\node[worker, right=1.6cm of q2] (w2) {Notification\\Worker};
\node[worker, right=1.6cm of q3] (w3) {Ranking\\Worker};

%% Scheduler
\node[box, below=1.5cm of api, fill=brand!15, minimum width=2.8cm] (sched)
      {Scheduler\\{\tiny cada 60 s}};

%% Flechas
\draw[arr] (api)  -- node[above,font=\tiny]{match.result\\loaded} (exch);
\draw[arr] (exch) -- (q1);
\draw[arr] (exch) -- (q2);
\draw[arr] (exch.south) |- (q3);
\draw[arr] (q1)   -- (w1);
\draw[arr] (q2)   -- (w2);
\draw[arr] (q3)   -- (w3);

%% Event chaining
\draw[arrg, dashed] (w1.north) -- ++(0,0.6) -| node[above,font=\tiny,pos=0.3]
      {scores.updated} (q3.north);

\draw[arr] (sched) -- node[left,font=\tiny]{round.closed} ++(2.5,0) |- (q3.south west);

\end{tikzpicture}
\caption{Flujo principal de eventos en Pick \& Serve.}
\label{fig:arch}
\end{figure}

La figura~\ref{fig:arch} resume la arquitectura. El backend publica el evento \texttt{match.result.loaded} al exchange de RabbitMQ; este lo enruta en \textit{fan-out} hacia tres colas independientes. El Scoring Worker calcula los puntos y, al terminar, publica a su vez \texttt{scores.updated} (encadenamiento de eventos), lo que dispara al Ranking Worker sin que exista ninguna dependencia directa entre ambos procesos. El Notification Worker atiende dos tipos de eventos. Un Scheduler independiente cierra rondas vencidas y publica \texttt{round.closed}.

% ─────────────────────────────────────────────────────────────
\section{Arquitectura Técnica}

\subsection{Stack tecnológico}

\begin{center}
\begin{tabular}{>{\bfseries}ll}
\toprule
Capa & Tecnología \\
\midrule
API & FastAPI 0.111 (Python) + Uvicorn \\
ORM / BD & SQLAlchemy 2.0 + PostgreSQL 16 \\
Broker & RabbitMQ 3 (Topic Exchange, colas durables) \\
Mensajería async & aio-pika 9.4 \\
Workers & Python (4 contenedores independientes) \\
Frontend & React 18 + TypeScript + Vite + nginx \\
Contenedores & Docker Compose v2 \\
\bottomrule
\end{tabular}
\end{center}

\subsection{Modelo de datos}

El esquema relacional cuenta con seis tablas: \texttt{users}, \texttt{tournaments}, \texttt{rounds}, \texttt{matches}, \texttt{predictions} y \texttt{rankings}. La tabla \texttt{notifications} almacena los mensajes in-app de cada usuario. Los campos clave del flujo son:

\begin{itemize}[noitemsep]
  \item \texttt{rounds.status} — \texttt{open} / \texttt{pending} / \texttt{closed}: controla en qué fase se aceptan predicciones.
  \item \texttt{matches.is\_final} — boolano que activa el bonus de puntos para la final del torneo.
  \item \texttt{predictions.points} — columna que persiste el resultado del scoring; los workers la verifican para garantizar idempotencia.
\end{itemize}

\subsection{Servicios Docker}

El sistema se levanta con un solo comando (\texttt{docker compose up}) e incluye ocho servicios:

\begin{center}
\small
\begin{tabular}{lll}
\toprule
Servicio & Puerto & Rol \\
\midrule
\texttt{postgres} & 5432 & Base de datos principal \\
\texttt{rabbitmq} & 5672 / 15672 & Broker + consola web \\
\texttt{backend} & 8000 & API REST (FastAPI) \\
\texttt{scoring-worker} & — & Calcula puntos \\
\texttt{notification-worker} & — & Crea notificaciones \\
\texttt{ranking-worker} & — & Actualiza ranking \\
\texttt{scheduler} & — & Cierra rondas por tiempo \\
\texttt{frontend} & 3000 & SPA servida por nginx \\
\bottomrule
\end{tabular}
\end{center}

% ─────────────────────────────────────────────────────────────
\section{Implementación de la Prueba de Concepto}

\subsection{Flujo principal: carga de un resultado}

El flujo de mayor interés para la evaluación de la arquitectura empieza cuando el administrador llama a \texttt{POST /admin/matches/\{id\}/result}. El endpoint registra el ganador en la base de datos, invoca \texttt{try\_advance\_round()} para determinar si todos los partidos de la ronda finalizaron y, de ser así, abrir automáticamente la siguiente, y publica el evento al exchange de RabbitMQ antes de devolver la respuesta HTTP. A partir de ahí, la API no hace nada más: el resto es responsabilidad de los workers.

\begin{lstlisting}[language=Python, caption={Publicación del evento tras cargar un resultado (admin.py).}]
await publish_event("match.result.loaded", {
    "match_id": match.id,
    "winner_player_id": winner_id,
    "round_id": match.round_id,
})
\end{lstlisting}

\subsection{Scoring Worker y encadenamiento de eventos}

El Scoring Worker consume \texttt{match.result.loaded}. Para cada predicción asociada al partido calcula los puntos según la tabla siguiente:

\begin{center}
\begin{tabular}{lc}
\toprule
Condición & Puntos \\
\midrule
Predicción incorrecta & 0 \\
Predicción correcta & +3 \\
Predicción correcta en final (\texttt{is\_final = true}) & +5 \\
\bottomrule
\end{tabular}
\end{center}

Si la predicción ya tiene puntos asignados (\texttt{predictions.points IS NOT NULL}), el worker la omite. Esto garantiza idempotencia ante reenvíos o fallas transitorias. Al concluir el procesamiento, publica \texttt{scores.updated}, lo que desencadena al Ranking Worker sin ningún acoplamiento directo entre ambos.

\subsection{Ranking Worker}

Responde a \texttt{scores.updated}. Suma los puntos acumulados por usuario, cuenta predicciones correctas y totales, ordena descendentemente y sobreescribe la tabla \texttt{rankings} con posiciones actualizadas. El frontend refresca este ranking cada cinco segundos.

\subsection{Notification Worker}

Atiende dos routing keys distintas:
\begin{itemize}[noitemsep]
  \item \texttt{match.result.loaded} → crea una notificación personalizada por usuario que haya hecho una predicción en ese partido (con ícono ✅ o ❌ según el resultado).
  \item \texttt{round.closed} → notifica a todos los usuarios no administradores que la ronda fue cerrada.
\end{itemize}

\subsection{Scheduler}

Corre en un contenedor separado con un loop de 60 segundos. Consulta las rondas con \texttt{status = 'open'} cuya fecha de inicio ya pasó y las cierra, publicando \texttt{round.closed} para que el Notification Worker actúe en consecuencia. Esto desacopla el tiempo del ciclo HTTP.

\subsection{Progresión de rondas (máquina de estados)}

\begin{figure}[H]
\centering
\begin{tikzpicture}[
  state/.style={ellipse, draw=brand, fill=brand!8, text centered,
                minimum width=2cm, font=\small},
  arr/.style={-Stealth, thick, color=brand}
]
\node[state] (pending) {pending};
\node[state, right=2.5cm of pending] (open)   {open};
\node[state, right=2.5cm of open]   (closed)  {closed};

\draw[arr] (pending) -- node[above,font=\tiny]{admin / scheduler} (open);
\draw[arr] (open)    -- node[above,font=\tiny]{todos los partidos\\finalizados} (closed);
\draw[arr, bend left=40] (closed.north) to node[above,font=\tiny]
      {abre la siguiente ronda} (open.north);
\end{tikzpicture}
\caption{Ciclo de vida de una ronda.}
\end{figure}

Cuando el último partido de una ronda abierta se marca como finalizado, el servicio \texttt{round\_progression.py} busca la siguiente ronda pendiente del mismo torneo (ordenada por \texttt{starts\_at}) y la abre. Esto permite que los usuarios comiencen a predecir en cuartos de final apenas termina la ronda de octavos, sin intervención manual del administrador.

% ─────────────────────────────────────────────────────────────
\section{Aspectos Destacados de la Implementación}

\subsection{Idempotencia en los workers}

Los workers están diseñados para ser seguros ante mensajes duplicados. La verificación de \texttt{predictions.points IS NOT NULL} en el Scoring Worker y la naturaleza \textit{upsert} del Ranking Worker garantizan que procesar el mismo evento dos veces no produce inconsistencias.

\subsection{Resiliencia en la conexión}

Tanto el backend como los workers implementan reintentos con backoff sobre la conexión a RabbitMQ y a PostgreSQL. El backend reintenta hasta 15 veces con intervalos de 5 segundos, permitiendo que el sistema se recupere de arranques lentos o reinicios en contenedores dependientes.

\subsection{Seed determinista}

Al arrancar, el backend ejecuta \texttt{seed.py} de forma idempotente: crea 10 usuarios, 2 torneos, 4 rondas y 6 partidos (incluyendo una final con bonus), con predicciones preexistentes para poder probar el flujo desde el primer momento sin configuración adicional.

\subsection{Testing}

La suite de pruebas cubre los cinco routers del backend (25 tests) usando SQLite en memoria y un mock de \texttt{publish\_event} que registra los eventos publicados. Esto permite verificar que cada endpoint dispara los eventos correctos sin necesidad de un broker real. El pipeline de CI en GitHub Actions ejecuta los tests de integración contra instancias reales de PostgreSQL y RabbitMQ.

% ─────────────────────────────────────────────────────────────
% ── ANEXOS ────────────────────────────────────────────────────
\newpage
\appendix
\section*{Anexo A — Capturas de la Consola RabbitMQ}
\addcontentsline{toc}{section}{Anexo A — Capturas de la Consola RabbitMQ}

Las siguientes imágenes muestran el estado del broker tras ejecutar el flujo de demostración descrito en el \texttt{README.md} del repositorio.

\begin{figure}[H]
  \centering
  % Reemplazar con la captura real del panel de exchanges
  \fbox{\parbox{0.85\textwidth}{\centering\vspace{2.5cm}
    \textit{[Captura: panel de Exchanges en localhost:15672 —\\
    exchange \texttt{pickandserve.events} de tipo topic, durable]}
  \vspace{2.5cm}}}
  \caption{Exchange \texttt{pickandserve.events} en la consola de administración de RabbitMQ.}
\end{figure}

\begin{figure}[H]
  \centering
  \fbox{\parbox{0.85\textwidth}{\centering\vspace{2.5cm}
    \textit{[Captura: lista de Queues —\\
    scoring.queue, notifications.queue, notifications.round.queue, ranking.queue\\
    con message rates y consumer count]}
  \vspace{2.5cm}}}
  \caption{Colas activas y sus consumidores durante la ejecución del PoC.}
\end{figure}

\begin{figure}[H]
  \centering
  \fbox{\parbox{0.85\textwidth}{\centering\vspace{2.5cm}
    \textit{[Captura: bindings del exchange \texttt{pickandserve.events} —\\
    routing keys: match.result.loaded, round.closed, scores.updated]}
  \vspace{2.5cm}}}
  \caption{Bindings del exchange, mostrando el enrutamiento por routing key.}
\end{figure}

\section*{Anexo B — Endpoints de la API}
\addcontentsline{toc}{section}{Anexo B — Endpoints de la API}

\begin{center}
\small
\begin{tabular}{lll}
\toprule
Método & Ruta & Descripción \\
\midrule
GET  & \texttt{/health}                           & Estado del servicio \\
GET  & \texttt{/auth/users}                       & Lista de usuarios \\
POST & \texttt{/auth/login}                       & Login por user\_id \\
GET  & \texttt{/rounds/open}                      & Rondas abiertas con partidos \\
POST & \texttt{/predictions}                      & Crear / actualizar predicción \\
GET  & \texttt{/predictions/me}                   & Predicciones del usuario \\
GET  & \texttt{/ranking}                          & Ranking global \\
GET  & \texttt{/notifications/me}                 & Notificaciones del usuario \\
PUT  & \texttt{/notifications/\{id\}/read}        & Marcar como leída \\
GET  & \texttt{/admin/matches}                    & Todos los partidos (admin) \\
POST & \texttt{/admin/matches/\{id\}/result}      & Cargar resultado (dispara evento) \\
POST & \texttt{/admin/rounds/\{id\}/close}        & Cerrar ronda manualmente \\
\bottomrule
\end{tabular}
\end{center}

\section*{Anexo C — Instrucciones de Ejecución}
\addcontentsline{toc}{section}{Anexo C — Instrucciones de Ejecución}

\begin{lstlisting}[language=bash, caption={Levantar el sistema completo.}]
git clone https://github.com/matiasmutz/pick-serve.git
cd pick-serve
docker compose up --build
\end{lstlisting}

Una vez que todos los servicios están en pie:

\begin{itemize}[noitemsep]
  \item Frontend: \url{http://localhost:3000}
  \item API (Swagger): \url{http://localhost:8000/docs}
  \item RabbitMQ Management: \url{http://localhost:15672} (guest / guest)
\end{itemize}

\textbf{Flujo de demostración sugerido:}
\begin{enumerate}[noitemsep]
  \item Ingresar con cualquier usuario jugador y realizar predicciones en los partidos de la ronda de Octavos del ATP Buenos Aires.
  \item Ingresar como administrador (\texttt{admin@pickserve.com}) y cargar el resultado de uno de los partidos.
  \item Observar en la consola de RabbitMQ cómo los mensajes fluyen por las colas.
  \item Volver al dashboard del jugador: las notificaciones y el ranking se actualizan sin necesidad de recargar la página.
  \item Cargar los resultados de los tres partidos de Octavos: el sistema cierra la ronda automáticamente y abre Cuartos de Final.
\end{enumerate}

\end{document}
